\begin{recipebox}{Coconut Chicken, Kumara \& Chickpea Curry}{3--4}{6--7 hrs (low) / 3--4 hrs (high)}

\begin{ingredients}
  \item 400\,g chicken breast
  \item 1 can coconut milk (400\,ml)
  \item 200\,g tinned tomatoes
  \item 2 tbsp red or green curry paste
  \item 100\,ml chicken stock
  \item 3 garlic cloves, minced
  \item 1 tsp fresh or ground ginger
  \item 1 tbsp fish sauce or soy sauce
  \item 1 tsp brown sugar or honey
  \item 1 medium kumara, peeled and diced into 2cm cubes
  \item 1 can chickpeas, drained and rinsed (400\,g)
  \item 100\,g frozen peas
  \item 200\,g jasmine or basmati rice (dry)
  \item 1 tbsp lime juice (optional)
\end{ingredients}

\begin{method}
  \item Mix the sauce: in the slow cooker, combine coconut milk, tinned tomatoes, curry paste, stock, garlic, ginger, fish sauce, and sugar. Stir to combine.
  \item Add chicken and kumara: place the whole chicken breast into the sauce. Add the diced kumara around the chicken --- it needs the full cook time to soften.
  \item Slow cook on LOW for 6--7 hours or HIGH for 3--4 hours, until the chicken is cooked through and the kumara is tender.
  \item Stir in the chickpeas. Cook for a further 30--60 minutes so they warm through and soak up some of the sauce.
  \item Remove the chicken breast and shred it with two forks. Return the shredded chicken to the slow cooker and stir through.
  \item Stir in the frozen peas. They'll be ready in about 10 minutes. Add lime juice if using, and taste for seasoning.
  \item Cook the rice according to packet instructions (usually \textasciitilde15 min). Fluff with a fork and serve with the curry.
\end{method}

\begin{notes}
Orange kumara works best here --- it's sweeter and pairs well with coconut. Gold kumara also works fine.

If you want even more body, mash a few chickpeas against the side of the slow cooker before serving --- it thickens the sauce nicely.

Ayam red curry paste is reliable and available at New World. Use 3 tbsp if you want more heat.
\end{notes}

\end{recipebox}